\section{Projectibility profiles}
\label{sec:7}

Section 1 made projectibility a secondary question. A positive access relation may support induction, but the comparative concept itself carries no predictions. This section records which grammatical profiles support further projection and what appears to stabilize them. The primary field is syntax, with targets such as agreement, syntactic-category licensing, and complement selection.

The distinction matters because \term{access transparency} and \term{projectibility profile} answer different questions. Applicability asks whether a case fills the schema's roles and controls. A relational verdict asks whether the relation reaches the bearer across the mediator. A profile verdict asks whether an observable base supports projection to further behaviour. The three verdicts should travel separately.

\subsection{Applicability and profile verdicts}\label{sec:7-1}

The umbrella schema is a comparative concept, not a single profile. Its grammatical applications vary in $X$, $P$, and $R$: QNs are the primary demonstration and TFRs a conditional grammatical illustration; Spanish and Tsez are controlled extensions; the Appendix~\ref{sec:B} cases are candidate topologies or boundary cases. No common stabilizer licenses projection across all of them. The schema's contribution is the four-gate applicability protocol and its three-way verdict, not a prediction about every case called \mention{transparent}.

Table 3 assigns the roles and controls case by case. A blank counterpart marks a missing control rather than a negative empirical result. The final column uses the terminology of the protocol: positive, non-instance, or indeterminate.

Table 3. Roles and controls filled, case by case.

{\def\LTcaptype{none}\footnotesize
\begin{longtable}[]{@{}
  >{\raggedright\arraybackslash}p{(\linewidth - 8\tabcolsep) * \real{0.17}}
  >{\raggedright\arraybackslash}p{(\linewidth - 8\tabcolsep) * \real{0.21}}
  >{\raggedright\arraybackslash}p{(\linewidth - 8\tabcolsep) * \real{0.17}}
  >{\raggedright\arraybackslash}p{(\linewidth - 8\tabcolsep) * \real{0.23}}
  >{\raggedright\arraybackslash}p{(\linewidth - 8\tabcolsep) * \real{0.22}}@{}}
\toprule\noalign{}
Configuration & Property, and its bearer & Mediator & Relation (diagnostic) & Concealing counterpart \\
\midrule\noalign{}
\endhead
\bottomrule\noalign{}
\endlastfoot
QN + \mention{of}-complement & number/count profile, borne by the complement & the QN, morphologically mismatched & agreement (verb form) & head-driven agreement \\
\addlinespace
TFR, parent & number profile, borne by the nucleus & \mention{what} & agreement (verb form) & ordinary fused relative \\
\addlinespace
TFR, attributional daughter & syntactic category, borne by the nucleus & \mention{what} & matrix selection (category flexibility) & evidential daughter, NP nuclei only \\
\addlinespace
Collective NP & member plurality, borne by the head itself & -- (same sign) & agreement (verb form) & -- \\
\addlinespace
Coordination & summed number, borne by no single conjunct & -- & agreement (verb form) & -- \\
\addlinespace
Deferred reference & number, borne by the designatum & the subject NP & agreement (verb form) & -- \\
\addlinespace
Spanish partitive subject & number, borne by the complement & \mention{mayoría} & agreement (verb form) & singular option, dispreferred \\
\addlinespace
Tsez complement clause & noun class, borne by the embedded absolutive & the clause & matrix agreement (class prefix) & class IV, the clausal default \\
\addlinespace
Aleut (App.~\ref{sec:B}) & person/number, borne by an unpronounced possessor & the possessed NP & anaphoric verb morphology plus subject case & overt-possessor alternant \\
\addlinespace
Lardil/Kayardild (App.~\ref{sec:B}) & clausal TAM, borne by the clause & the dependency domain & dependent-nominal morphology & -- \\
\addlinespace
Chamorro (App.~\ref{sec:B}) & grammatical relation of the extractee & the extraction path & path morphology & -- \\
\addlinespace
Obviation (App.~\ref{sec:B}) & possessor obviation, borne by the possessor & the possessed NP & predicate tracks the head's derived feature & -- \\
\end{longtable}
}

Table 4. Applicability, relational, and profile verdicts.

{\def\LTcaptype{none}
\begin{longtable}[]{@{}
  >{\raggedright\arraybackslash}p{(\linewidth - 4\tabcolsep) * \real{0.25}}
  >{\raggedright\arraybackslash}p{(\linewidth - 4\tabcolsep) * \real{0.39}}
  >{\raggedright\arraybackslash}p{(\linewidth - 4\tabcolsep) * \real{0.36}}@{}}
\toprule\noalign{}
Case & Applicability and relational verdict & Profile verdict \\
\midrule\noalign{}
\endhead
\bottomrule\noalign{}
\endlastfoot
Number-transparent QNs (§\ref{sec:2-1}) & Positive: complement number/count reaches agreement through the QN & Strongest candidate; selected agreement and determiner tests support projection, but the full cluster remains open \\
\addlinespace
Transparent free relatives (§\ref{sec:2-5}) & Positive conditional on Kim's constructional analysis: nucleus properties reach agreement and selection through \mention{what} & Second syntactic candidate; pilot verb-class evidence, not a completed profile audit \\
\addlinespace
Collective NPs (§\ref{sec:2-2}) & Non-instance: bearer and mediator are layers of one sign, not distinct elements & None claimed \\
\addlinespace
Ordinary coordination (§\ref{sec:2-3}) & Non-instance: no mediator conceals the summed features & None \\
\addlinespace
Deferred reference (§\ref{sec:2-3}) & Non-instance: no bearer occurs within $X$ & None \\
\addlinespace
Spanish partitive subjects (§\ref{sec:6-topologies}) & Positive with a weak counterpart: complement number reaches agreement despite a dispreferred singular option & No English-like profile shown \\
\addlinespace
Tsez long-distance agreement (§\ref{sec:6-topologies}) & Positive: embedded noun class reaches matrix agreement against a class-IV counterpart & None claimed \\
\addlinespace
Aleut (§\ref{sec:B}) & Positive with qualification under Merchant's analysis of a null possessor & None claimed \\
\addlinespace
Lardil/Kayardild (§\ref{sec:B}) & Indeterminate: the inward relation is described, but a same-family counterpart hasn't been established & None claimed \\
\addlinespace
Chamorro (§\ref{sec:B}) & Indeterminate: path morphology is a candidate relation, without the required counterpart & None claimed \\
\addlinespace
Obviation (§\ref{sec:B}) & Non-instance or boundary case: the predicate tracks a feature relayed onto the head & None claimed \\
\end{longtable}
}

The table separates a missing control from a failed relation. A case can be positive while its profile remains untested, as with Tsez and the Spanish partitive. A proposed profile can also weaken while the local grammatical relation survives, as the sparse inanimate \mention{bunch} result shows for the QN discussion. Indeterminate cases shouldn't be promoted to positive instances merely because their topology is suggestive.

\subsection{What the primary profiles project}\label{sec:7-2}

The QN class is the strongest candidate because its observable base is independently available. Known members form a small conventional inventory, while candidate members can be admitted on quantificational semantics and frequent \mention{of}-complement distribution before their projected behaviours are examined. Determiner frames, complement selection, and agreement direction covary by subgroup. The COCA audit supports selected directions, but asymmetric filtering and sparse subject-position data leave the end-to-end cluster open.

The \mention{piles} test illustrates the proper evidential grade. Its determiner prediction received partial prospective support and revealed a broader literal-versus-quantificational sense split across the plural-only subgroup. Its agreement prediction wasn't testable because subject-position tokens were too sparse. The result supports candidate status and identifies the next feasible method, probably acceptability judgments, rather than validating the whole profile.

The status of \mention{rest} is an analytical fork. Under a quantity-noun analysis, its determiner and agreement behaviour can enter the candidate-member comparison. Under a part/subset analysis, it remains a useful access-transparency case but shouldn't supply evidence for the quantity-QN discovery criteria. Keeping those questions apart protects the profile claim from a lexical classification dispute.

TFRs are a second syntactic candidate. Their parent/daughter split and verb-class restriction suggest a cluster stabilized by constructional inheritance. The current corpus check supports the attributional/evidential contrast, but the result remains conditional on Kim's constructional analysis and lacks the category-discriminating judgments needed for a completed profile. The Spanish and Tsez cases show that the access relation travels; they don't show that the English QN cluster travels with it.

The form--meaning cases in §\ref{sec:3} are boundary applications. Their relations and diagnostics may support field-specific profiles, but this paper makes no claim about a shared compound, constructional, and processing cluster.

\subsection{Stabilizers and diagnostics}\label{sec:7-3}

\term{Lexical convention} is the proposed stabilizer for the QN profile. Speakers acquire item-specific slots such as \mention{a lot of} X, \mention{plenty of} X, \mention{the rest of} X, and \mention{a number of} X, together with their agreement and complement-selection behaviour. This account explains why the class is small and why boundary members and sense splits show graded behaviour. It doesn't require Boyd-style homeostasis. The broader countability analysis in \autocite{reynolds2026countability} supplies a possible causal-cluster connection, but the QN argument here doesn't depend on it.

\term{Constructional inheritance} is the proposed stabilizer for the TFR profile. The parent and attributional daughter support different access conditions, and the relevant candidate projections follow the verb-class hierarchy. That proposal is stronger than the current evidence warrants as a completed result, but it identifies a concrete source of covariance for the joint paper to test.

Slater's \autocite*{slater-2015-natural-kindness} stable property cluster (SPC) account asks whether a cluster is stably coinstantiated enough to support the inferential uses assigned to it. The QN profile is the strongest candidate for this diagnostic because determiner, complement, and agreement behaviours covary by subgroup. The TFR profile is a weaker candidate because its verb-class evidence is still pilot-level. Khalidi's \autocite*{khalidi-2018-causal-networks} causal-network account asks a different question about ordered causal stabilizers. The QN class appears conventionally stabilized rather than causally hierarchical, so SPC is the relevant assessment here; no Khalidian result is claimed for the umbrella.

Kindhood is downstream. A profile that supports stable projection may earn a local-kind label; a profile with partial support remains a candidate; a useful descriptive category may remain non-projectible. The comparative concept doesn't settle that question in advance.

\subsection{The resulting claim}\label{sec:7-4}

\term{Access transparency} is a relational comparative concept, not an entity or a natural kind. Its schema identifies a bearer, mediator, relation, diagnostic, and concealing counterpart. The four-gate protocol distinguishes genuine mediation from feature summation, external reference, and analysis-dependent cases that should remain indeterminate.

Three levels stay distinct:

\begin{itemize}
\tightlist
\item
  \term{Comparative concept}: access transparency, the cross-linguistic methodological tool defined by the schema of §\ref{sec:1}.
\item
  \term{Descriptive category}: language-particular grammatical content, such as English number-transparent QNs, Spanish partitive subjects, and Tsez long-distance agreement.
\item
  \term{Projectibility profile}: an observable base that licenses non-trivial projection to a target, stabilized by an identifiable source, with local-kind status as a downstream verdict.
\end{itemize}

The contribution is to apply projectibility-first discipline to a relational concept without treating every filled schema as a projectible kind. Applicability comes first. Profile evidence comes second. The distinction lets the paper preserve a useful comparative concept while remaining cautious about which language-particular patterns earn further inference.
